\section{相と溶液から分離操作へ}
蒸留は、加熱して一部を気体にすれば組成が変わることを利用する操作である。その理解には「液体は沸点でしか気化しない」という直感を捨て、気相と液相の平衡、溶媒に他成分を加えた効果、両相の組成差を順に考える必要がある。熱力学の相平衡式は第\ref{theory:thermo:phase}節で導いた。ここでは容器の条件と分離の技術的意味を結び付ける。

\subsection{化学ポテンシャルが等しいとは何か}
\label{theory:phase_sep:chemical_potential}
純物質では化学ポテンシャル \(\mu=G/n\) は1 molをその相へ加えたときのギブズエネルギーの増分である。同じ温度・圧力の液体から蒸気へ微小量 \(dn\) を移すと
\[
 dG=(\mu_g-\mu_l)\,dn
\]
となる。両相が平衡なら、どちらへの移動でも自由エネルギーを減らせず \(\mu_g=\mu_l\) である。固液、固気でも同じ条件を用いる。これは分子の移動が止まることではなく、両方向の正味流束が等しいことを表す。

\subsection{定温の飽和蒸気圧と、液体が残る条件}
\label{theory:phase_sep:vapor}
純物質の気液二相が平衡にあると、温度を指定すれば圧力は飽和蒸気圧 \(p_{\mathrm{sat}}(T)\) に決まる。容器の体積を増すと一時的に気相の密度が下がるが、液体が追加で蒸発して平衡圧力を回復する。温度が同じなら、気相の量は増えても圧力は同じである。

この議論は、膨張後にも液体が残る場合に限る。液体が全部蒸発すれば二相の平衡条件は失われ、単相気体の状態方程式と全物質量で圧力を決める。「体積を2倍にし十分静置した」と指定しても、液体の初期量が少なければ残存は保証されない。

\subsection{蒸発と沸騰を分ける}
\label{theory:phase_sep:boiling}
沸点より低い温度でも表面の分子は液体から気体へ移る。乾燥した空気の中で水が減ることは、沸騰せずに蒸発する身近な例である。沸騰は液体内部で蒸気泡が成長する条件であり、巨視的には蒸気圧が周囲の圧力に達する温度を沸点とする。微小気泡では曲率による圧力差も関係するが、通常の試験問題では省略する。

純液体を一定外圧で加熱し、二相が共存する間は温度が沸点に保たれ、熱は蒸発潜熱へ使われる。多成分液体では蒸発によって組成が変わるため、一定外圧でも沸点が変化し得る。液体が全部なくなった後の蒸気の昇温も別である。

\subsection{密封容器の全圧と水蒸気の分圧}
\label{theory:phase_sep:sealed}
室温の空気と水を密封した剛体容器では、近似的に
\[
 p_{\mathrm{total}}=p_{\mathrm{air}}+p_{\mathrm{water}},
 \qquad p_{\mathrm{air}}\simeq\frac{n_{\mathrm{air}}RT}{V_g}
\]
となる。水が残れば \(p_{\mathrm{water}}\) は飽和蒸気圧に対応するが、それが容器の全圧ではない。100°Cで水の飽和蒸気圧が約1気圧でも、空気の分圧が加わるため全圧は1気圧を超える。開放容器の沸点をそのまま使えない理由である。空気を除いた容器や気相部のない充填では条件が変わる。

\subsection{潜熱、エントロピー、蒸気圧の温度変化}
\label{theory:phase_sep:latent}
同じ温度・圧力の気液平衡では \(\Delta G_{\mathrm{vap}}=0\) なので
\[
 \Delta S_{\mathrm{vap}}=\frac{\Delta H_{\mathrm{vap}}}{T}.
\]
蒸発と凝縮は同じ端点を逆向きに結ぶためエンタルピー差は等大逆符号である。蒸気を理想気体とし液体体積を無視すると、クラペイロン式は
\[
 \frac{d\ln p_{\mathrm{sat}}}{dT}=\frac{\Delta H_{\mathrm{vap}}}{RT^2}
\]
へ近似できる。潜熱一定はこの微分式から積分形へ移るための追加近似であり、臨界点まで一定とすることはできない\parencite{mit_clapeyron2008}。

\subsection{不揮発性溶質はなぜ蒸気圧を下げるか}
\label{theory:phase_sep:raoult}
溶媒をAとし、理想溶液では液相の化学ポテンシャルが
\(\mu_A=\mu_A^*+RT\ln x_A\) となる。\(x_A<1\) なので純溶媒より低い。気相の化学ポテンシャルと等置すると、ラウールの法則 \(p_A=x_Ap_A^*\) が得られる。不揮発性溶質Bは蒸気にほとんど寄与しないため、同じ温度で溶液の蒸気圧は純溶媒より低くなる。実在溶液ではモル分率を活量に置き換える\parencite{mit_solutions2008}。

\subsection{凝固点降下：濃度でなく粒子数を数える}
\label{theory:phase_sep:freezing}
溶媒の液相の化学ポテンシャルが混合で低下すると、固相との一致に必要な温度も変わる。溶質が固相へ入らない希薄溶液では、融解平衡式を純溶媒の凝固点の近くで展開して
\[
 \Delta T_f=K_f m_{\mathrm{particles}}=iK_fm_{\mathrm{formula}},
 \qquad
 K_f=\frac{RT_f^{*\,2}M_A}{\Delta H_{\mathrm{fus},m}}
\]
を得る。\(M_A\) は溶媒のモル質量をkg/molで表した値、\(m\) は溶媒1 kg当たりの物質量である。\(i\) は式量単位に対する溶質粒子数の倍率であり、完全解離するNaClなら2となる。正の量 \(\Delta T_f\) は降下幅なので、凝固温度は \(T_f^*-\Delta T_f\) である。実際の電解質では会合と相互作用により単純な粒子数の式からずれる\parencite{mit_colligative2008}。

\subsection{三重点、臨界点、過冷却}
\label{theory:phase_sep:critical}
三重点は一成分の固体・液体・気体が同時に平衡共存する点である。臨界点は気液共存曲線の終端であり、両相の密度差が消え、巨視的な界面がなくなる。臨界温度より高い状態で圧力を増して高密度流体にしても、通常の気液相転移を経て液化する過程ではない。

\label{theory:phase_sep:isochore}
剛体密閉容器では全質量 \(M\) と容積 \(V\) が固定され、平均密度 \(\bar\rho=M/V\) が一定である。二相共存の密度を \(\rho_g(T),\rho_l(T)\)、液相の体積割合を \(\phi\) とすると
\[
 \bar\rho=(1-\phi)\rho_g+\phi\rho_l,\qquad0\le\phi\le1.
\]
水の臨界密度は \(322\,\mathrm{kg\,m^{-3}}\)、参照臨界温度は647.096 Kである\parencite{iapws95_2018}。平均密度を臨界密度に選べば共存状態から臨界へ近づける。一方、別の平均密度なら臨界に達する前にどちらかの相が消える場合がある。「密閉して加熱すれば必ず臨界へ達する」と一般化しない。

\subsection{水の準安定性と温度尺度}
\label{theory:phase_sep:metastable}
\label{theory:phase_sep:water}
平衡で氷が安定でも、氷核の生成が遅ければ液体の水は0°C未満まで残る。これが過冷却であり、安定相の選択と相転移速度を分ける必要がある。IAPWSの過冷却水の指針は、大気圧近傍を含め、融解温度と均一核生成温度の間の準安定液体を扱う\parencite{iapws_supercooled2015}。水の三重点の0.01°CはITS-90の慣用値として読む。現行SIの温度定義との違いは第\ref{theory:units:temperature}節を参照する。

\subsection{実在気体と熱容量の近似}
\label{theory:phase_sep:realgas}
\label{theory:phase_sep:capacity}
ファンデルワールス式 \((p+a/V_m^2)(V_m-b)=RT\) の \(a\) は引力、\(b\) は分子の有限サイズによる排除体積を表す。CO2と窒素の理想性は同じ温度・圧力で比較する。標準的な低圧条件ではCO2の方が非理想性が大きいが、全状態で順位が固定するという意味ではない。

定圧と定積の熱量の差を膨張仕事だけで説明できるのは \(U=U(T)\) の理想気体である。実在気体では同じ温度でも最終体積による内部エネルギー差を持つ。熱容量差の一般式は第\ref{theory:thermo:solidcapacity}節に示した。

\subsection{自由膨張とギブズエネルギーの条件}
\label{theory:phase_sep:free}
\label{theory:phase_sep:gibbs}
真空へ自由膨張する断熱閉鎖系では熱も外部への体積仕事も0なので内部エネルギーは一定である。理想気体の温度は変わらない。一般の単相流体では
\[
 dU=C_VdT+\left[T\left(\frac{\partial p}{\partial T}\right)_V-p\right]dV
\]
より、自由膨張の温度変化は状態方程式に依存する。理想気体は角括弧が0、ファンデルワールス模型は \(a n^2/V^2>0\) となる。断熱膨張一般の冷却則と取り違えない。

\(G=U+pV-TS\) は温度・圧力を自然変数とする自由エネルギーである。「定圧条件の自由エネルギー」という略称だけから、任意の定圧過程でGが減少すると推論しない。自発変化の判定は温度も一定などの条件を要する\parencite{iupac_greenbook2012}。

\subsection{天秤と空気の浮力}
\label{theory:phase_sep:buoyancy}
空気中の試料の見掛けの重量は、空気密度 \(\rho_{\mathrm{air}}\)、試料体積 \(V_s\)、質量 \(M_s\) に対し \((M_s-\rho_{\mathrm{air}}V_s)g\) である。分銅側にも浮力が働くが、両側の体積が違えば完全には打ち消されない。温度・湿度で空気密度が変わるため精密な質量測定に補正が必要となる。湿度上昇が密度へ与える効果は同じ温度・全圧で比較する。

\subsection{蒸留を定量化する気液平衡比}
\label{theory:phase_sep:distillation}
液相の成分iのモル分率を \(x_i\)、気相を \(y_i\) とすると、気液平衡比と2成分の相対揮発度は
\[
 K_i=\frac{y_i}{x_i},\qquad
 \alpha_{ij}=\frac{K_i}{K_j}
 =\frac{y_i/x_i}{y_j/x_j}.
\]
気相へ移りやすい成分が蒸気へ濃縮されれば、蒸気を取り出して凝縮することで組成を変えられる。相対揮発度が1に近いほど分離は難しい。一定の相対揮発度を用いる二成分近似では \(y=\alpha x/[1+(\alpha-1)x]\) となる\parencite{nptel_distillation1}。

このK値は相分配を表し、化学反応の平衡定数とは定義が違う。「平衡定数」が蒸留と最も関係が薄いという用語問題では、どちらを指すか原典に明記されない。この曖昧さを消して「蒸留は平衡と無関係」と教えてはいけない。

\subsection{理論段数、還流、充填塔}
\label{theory:phase_sep:stages}
1段の出口気液が平衡に達する理想化を理論段と呼び、目的の分離に必要な段数を理論段数とする。塔頂の凝縮液の一部を戻す還流によって、下向き液体と上向き蒸気を向流接触させ、蒸発と部分凝縮を反復する。棚段塔は段ごとの接触、充填塔は充填物の表面での連続した接触を利用する。段数や還流は同じ気液平衡のもとで接触効果を高める量である。

\subsection{共沸は通常精留の限界となる}
\label{theory:phase_sep:azeotrope}
共沸組成では気液組成が同じになり、蒸発させても残液の組成が変わらない。その圧力での通常精留では完全分離できない\parencite{nptel_azeotrope}。還流や段数を増しても、気液平衡曲線の組成差が0という制約を消せない。圧力変更、第三成分の利用、別の分離方法などは、元の通常精留と異なる操作である。試験のエタノール「96\%」には濃度基準が書かれないため、質量・体積・モルパーセントを混同しない。

\subsection{水蒸気蒸留と減圧蒸留}
\label{theory:phase_sep:steam}
ほぼ混和しない2液が共存すると、それぞれの蒸気分圧が全圧へ寄与する。水蒸気蒸留はこの条件で対象の揮発成分を水蒸気とともに運ぶ。減圧蒸留は外圧を下げて低い温度で沸騰させ、熱分解を起こしやすい物質の分離に用いる。両者の目的と圧力・混和性の条件は異なる。

\subsection{低温面へ気体を捕捉する真空ポンプ}
\label{theory:phase_sep:vacuum}
クライオポンプは冷却面で気体を凝縮し、凝縮しにくい気体を低温吸着材へ捕捉する。気体を系内へ蓄えるため捕捉容量を持ち、加熱して再生する必要がある\parencite{leybold_cryopump}。ターボ分子・拡散ポンプの運動量輸送、ダイヤフラム・ルーツ型の機械的移送とは排気対象の扱いが違う。

\subsection{コロイドは何が分散しているかで読む}
\label{theory:phase_sep:colloid}
普通の溶液は分子やイオンが均一に混じった系である。コロイドでは、より大きい粒子や液滴などが分散媒に分散する。液滴が別の液体へ分散した乳濁液がエマルションであり、マヨネーズはその例である。分散粒子の光散乱で光路が見える現象をチンダル現象という\parencite{openstax_colloids}。

\label{theory:phase_sep:micelle}
石けん分子は親水部と疎水部を持ち、温度・濃度などの条件が適すると会合してミセルを作る。ミセルと、石けんが表面へ並んで安定化した油滴とは区別する。臨界ミセル濃度より薄い系まで同じ粒子状態とはしない。「コロイド溶液」の慣用語があっても、溶解という語だけでは分散種や尺度を特定できない。

\section{この単元の問題を読む順序}
まず純物質か混合物かを決め、温度・圧力・体積・物質量のどれが固定されるかを確認する。次に相が残るか、安定平衡か準安定かを区別し、気相の組成と量を読む。最後に分離装置の工夫が接触効率を変えるのか、気液平衡自体を変えるのかを判断する。以下の初稿で残した条件不足は原典の省略であり、公式正答の一致によって解消したとは扱わない。
